\documentclass[11pt]{article}
\usepackage[T1]{fontenc}\usepackage{lmodern}
\usepackage{amsmath,amssymb,amsthm,booktabs,hyperref,geometry,longtable}
\geometry{margin=2.3cm}
\newtheorem{theorem}{Theorem}
\title{Polynomial continued fractions with factorial reduction:\\ hidden low-degree equivalences, a staged identification pipeline,\\ and eight unidentified candidates}
\author{Vichien Fugsukjit\\ Independent researcher, Bangkok, Thailand\\ \texttt{9chienvi9@gmail.com}}
\date{DRAFT v2, \today}
\begin{document}\maketitle
\begin{abstract}
We search polynomial continued fractions (PCFs) with quadratic partial denominators and quartic partial numerators for full
factorial reduction (FR) and remove every case explained by known mechanisms (the family of Conjecture~1.3 of Ben David et al.,
deflation, rational or quadratic limits, M\"obius transforms of standard constants and of H.~Cohen's databases, and Euler-type
PCFs detected through hypergeometric solutions of the recurrence). One PCF that passed all numerical tests is shown to be of Euler
type in disguise (Theorem~1). Fifteen PCFs remained whose order-four operator is irreducible over $\mathbb Q(z)$ and whose limits
match none of 21{,}984 constants at 1{,}000 digits. A further search for M\"obius equivalences with PCFs of LOWER degree shows that
seven of the fifteen are disguised degree-$(1,2)$ PCFs, i.e.\ ratios of contiguous Gauss hypergeometric functions (agreement to
170 digits); for three of them we give exact proofs by polynomial coboundaries (Theorem~2), revealing a simple hiding mechanism
that standard deflation tests do not detect. The remaining eight passed every test we could devise (Gamma products at 500 digits,
4{,}310 L-values of rational newforms, low-degree searches over about $10^9$ partner PCFs, monodromy, parameter deformation)
and are published as candidates, not as new constants. We release the data, the code and a staged identification-and-proof
pipeline that would have identified the seven on the first day.
\end{abstract}

\section{Conventions}
A PCF is $x=a(0)+\cfrac{b(1)}{a(1)+\cfrac{b(2)}{a(2)+\cdots}}$ with $a,b\in\mathbb Z[n]$; $P_n,Q_n$ satisfy $u_n=a(n)u_{n-1}+b(n)u_{n-2}$
with $P_{-1}=1$, $P_0=a(0)$, $Q_{-1}=0$, $Q_0=1$. In matrix form $(u_{n-1},u_n)=(u_{n-2},u_{n-1})M(n)$ with $M(n)=\left(\begin{smallmatrix}0&b(n)\\1&a(n)\end{smallmatrix}\right)$.
Full FR means that $S_N=\frac1N\log(|Q_N|/\gcd(P_N,Q_N))$ stays bounded; numerically $|S_{800}-S_{200}|<0.15$ (exact integers).

\section{Search and filters}
\begin{enumerate}
\item Brute force over about $2.3\times10^9$ PCFs ($\deg a=2$, $b$ a product of four linear factors) with M\"obius hashing against
standard constants: 239 unidentified FR-candidates, 45 with full FR, 23 outside Conjecture~1.3 and a ${}_3F_2$ scan.
\item Extended search on the six $a(n)$ of these 23 with $b(n)=s\prod_{i=1}^4(k_in+t_i)$: 5,274 PCFs with full FR.
\item Removed: Conjecture~1.3 incl.\ index shifts (846), slow convergence (15), rational or quadratic limits (313), deflation (3,684);
M\"obius clustering left 347 new clusters; a Cohen-database limit (1), $\pi\sqrt3$ (1), ${}_3F_2$ (2).
\item Euler type (order-one right factor of the recurrence over $\mathbb Q(n)$, ore\_algebra): 61 + 1 cases (Section~3).
\item $L=\theta^2(\theta-1)^2-z\,a(\theta+1)\theta^2-z^2\,b(\theta+2)$ annihilates $\sum_nQ_nz^n/(n!)^2$; fifteen operators are
irreducible over $\mathbb Q(z)$ (labels C0--C14, SHA-256 time-stamped before disclosure).
\end{enumerate}

\section{A hidden Euler-type example}
\begin{theorem}
For $a(n)=4n^2+10n-3$, $b(n)=-n^2(n+4)(3n-4)$: $P_n=-(n+1)!\prod_{m=0}^{n-1}(3m+2)\,(n^2+3n+3)$ and $x=1/(S-1/3)$ with
$S=-\frac1{24}\sum_{k\ge1}\frac{k!\,(k+2)(k+3)(k+4)}{(2/3)_k\,3^k\,(k^2+k+1)(k^2+3k+3)}$.
\end{theorem}
\begin{proof}[Proof sketch]
$r(n)=P_n/P_{n-1}=(n+1)(3n-1)(n^2+3n+3)/(n^2+n+1)$ satisfies $a(n)r(n-1)+b(n)=r(n)r(n-1)$ in $\mathbb Q(n)$; the Casoratian of
$(P,Q)$ gives the series. Checked to 120 digits.
\end{proof}
With the same $a(n)$, $b(n)=(n+3)(2n-5)(2n-1)(3n-1)$ and $(n+1)(2n-3)(2n+1)(3n-5)$ give M\"obius transforms of
$\Gamma(1/3)^3\sqrt3/\pi$-type periods (250 digits); both operators factor as $[1,3]$.

\section{Seven of the fifteen are disguised \texorpdfstring{degree-$(1,2)$}{degree-(1,2)} PCFs}
For each candidate $x$ we searched PCFs $y$ of lower degree (Batches 12--18: $\deg a\le2$, $\deg b\le2$, coefficient ranges up to
$\pm60$ for $a$ and up to $\pm40$ or factored with $k\le12,|t|\le30$ for $b$; about $10^9$ partners) for a M\"obius relation
$y=(p+qx)/(r+sx)$, $|p|,\dots,|s|\le8$, with planted, natural and random controls in every run. Every match was recomputed with
exact recurrences to 160 digits.
\begin{center}{\small\begin{tabular}{llllc}\toprule
cand. & partner $a(n)$ & partner $b(n)$ & relation & status\\\midrule
C0 & $2n+3$ & $(3n+5)(5n+1)$ & $y=3-x$ & proven\\
C10 & $n-3$ & $(2n+1)(3n+2)$ & $y=(8-6x)/(4-5x)$ & proven\\
C11 & $n-3$ & $(3n-1)(4n+1)$ & $y=(8-3x)/(4+x)$ & proven\\
C1 & $2n-1$ & $(3n-5)(5n-21)$ & $y=-(5+4x)/5$ & 171 digits\\
C5 & $6n-3$ & $9(3n-2)(5n-2)$ & $y=(x-3)/(-4-x)$ & 171 digits\\
C7 & $2n$ & $4(2n+3)(3n+1)$ & $y=-3x/(5-x)$ & 171 digits\\
C14 & $2n+8$ & $4(4n+1)(5n-7)$ & $y=(4+x)/6$ & 172 digits\\\bottomrule
\end{tabular}}\end{center}
Degree-$(1,2)$ PCFs are ratios of contiguous ${}_2F_1$ functions; hence none of the seven is a new constant.
\begin{theorem}
Let $M_1$ be the matrix of candidate C0 ($a=4n^2+10n-3$, $b=(2n-3)(2n+1)(3n+2)(5n-4)$) and $M_2$ that of its partner. Then
$M_1(n)\,U(n)=(2n+1)\,U(n-1)\,M_2(n-1)$ with
\[U(n)=\begin{pmatrix}\frac{10}{3}n^2+\frac73n-2 & -10n^3-\frac{41}{3}n^2+6n+\frac53\\[2pt] -\frac13 & n\end{pmatrix},\]
$\det U(n)=-\frac59(2n-1)(2n+1)\neq0$ for integers $n$. Analogous identities hold for C10 ($\lambda=3n+4$) and C11 ($\lambda=3n+1$).
\end{theorem}
The identities are polynomial and were verified symbolically; they make the convergent matrices equivalent and imply the M\"obius
relations of the table. Mechanism: the candidate is its degree-$(1,2)$ partner shifted by one step and wrapped by one linear factor
$\lambda(n)$ of its own $b(n)$, mixed by a cubic polynomial gauge $U(n)$; the partner's $b$ is then a product of two linear factors
of the candidate's $b$, which turns the search into a cheap filter. C1, C5, C7 and C14 are not obtained this way within the searched
gauge degrees (up to 5), nor by a two-step (contraction) ansatz whose positive control failed, so their mechanism is open.

\section{The eight unidentified candidates}
\begin{center}{\small\begin{tabular}{lllll}\toprule
 & $a(n)$ & $b(n)$ & $x$ & $\rho$ \\\midrule
C2 & $3n^2+8n+8$ & $2(2n+3)(3n+1)(3n+2)(3n+8)$ & $33.5147685956496535\ldots$ & $-187/42$\\
C3 & $3n^2+8n+8$ & $2(2n+3)(3n-1)(3n+8)(5n+3)$ & $27.7112473245165813\ldots$ & $-187/45$\\
C4 & $3n^2-2n+2$ & $2(2n-1)(3n-8)(3n-2)(5n-8)$ & $-9.6178657191277317\ldots$ & $3/5$\\
C6 & $4n^2+9n+1$ & $(2n+3)(4n+1)(4n+9)(6n+7)$ & $39.4516159585089660\ldots$ & $-123/28$\\
C8 & $6n^2+13n+5$ & $2(2n+1)(2n+5)(3n+4)(3n+7)$ & $36.8847177941179069\ldots$ & $-83/18$\\
C9 & $6n^2+13n+5$ & $(3n-10)(3n-4)(3n+10)(5n+3)$ & $-41.3047171139817542\ldots$ & $-121/60$\\
C12 & $4n^2+10n-3$ & $(2n-5)(2n+3)(3n+5)(5n-9)$ & $49.6956732085298041\ldots$ & $-39/20$\\
C13 & $6n^2+13n+5$ & $2(2n+5)(2n+9)(4n+3)(5n-8)$ & $-54.0753445873853959\ldots$ & $-241/52$\\\bottomrule
\end{tabular}}\end{center}
$\rho$: non-integer local exponent of $L$ at its dominant singularity.

{\sloppy\paragraph{Tests passed by all eight (none identified).} Not in Conjecture~1.3; not deflatable; limit not rational or quadratic;
not M\"obius ($|c|\le10^5$) of 36 standard constants, of the 635 limits of arXiv:2409.06086v4, of evaluable limits of arXiv:2607.06581,
or of the 171 euler2ai limits; no hypergeometric solution; operator irreducible; PSLQ $[1,x,c,xc]$ against 21,984 constants at 1,000
digits; no match in the OEIS; connection coefficients not $\Gamma$-products at 150 and 480 digits (planted controls recovered,
incl.\ a 53-element basis); no M\"obius relation ($|c|\le10^8$) with 4,310 L-values of rational newforms of weight 2--4 and level
$\le100$ (300 digits); no low-degree M\"obius partner in the searches of Section~4; traces of short words in the numerical monodromy
are not algebraic of degree $\le4$ (controls: hypergeometric operators give 100\% algebraic traces); shifting one linear factor of
$b$ by its natural step gives no M\"obius relation between neighbours of degree $\le4$.\par}

\section{Limitations}
Novelty is not established: unpublished databases (e.g.\ the Ramanujan Library) were not available to us. C12 was not completed
by the widest mechanism filter. The identifications of C1, C5, C7, C14 are numerical (170 digits). Irreducibility is over
$\mathbb Q(z)$ only. Full FR is a numerical criterion. The LMFDB has no search for decimal constants.

\section*{Data and code}
All plans were time-stamped on GitHub (github.com/vichienFF/ramanujan-explorer) before the corresponding runs; scripts, logs,
results and the pipeline \texttt{identify\_pipeline.py} (stages: basic constants, known structures, mechanism filter, low-degree
search, constant libraries, operator analysis, coboundary proof; self-test included) are released with this preprint.
\section*{Acknowledgements}
Computations used SageMath, ore\_algebra, ramanujantools, mpmath, PARI/GP and Kaggle notebooks. The author used AI assistants
(Anthropic Claude for code and analysis, OpenAI GPT as an independent checker); all claims are reproducible with the released code.

\begin{thebibliography}{11}
\bibitem{RM} G.~Raayoni et al., Generating conjectures on fundamental constants with the Ramanujan Machine, Nature 590 (2021) 67--73.
\bibitem{PNAS} R.~Elimelech et al., Algorithm-assisted discovery of an intrinsic order among mathematical constants, PNAS 121 (2024), doi:10.1073/pnas.2321440121.
\bibitem{BD} N.~Ben David et al., On the connection between irrationality measures and polynomial continued fractions, arXiv:2111.04468.
\bibitem{OD} O.~David, On Euler polynomial continued fractions, arXiv:2308.02567.
\bibitem{CMF} O.~David et al., The conservative matrix field, arXiv:2303.09318.
\bibitem{E2AI} T.~Raz et al., From Euler to AI: unifying formulas for mathematical constants, arXiv:2502.17533.
\bibitem{C1} H.~Cohen, A database of continued fractions of polynomial type, arXiv:2409.06086.
\bibitem{C2} H.~Cohen, Continued fractions of polynomial type: theory and encyclopedic dictionary, arXiv:2607.06581.
\bibitem{Hyper} M.~Petkov\v sek, Hypergeometric solutions of linear recurrences with polynomial coefficients, J.\ Symbolic Comput.\ 14 (1992) 243--264.
\bibitem{ore} M.~Kauers, M.~Jaroschek, F.~Johansson, Ore polynomials in Sage, LNCS 8942 (2015).
\bibitem{Apery} R.~Ap\'ery, Irrationalit\'e de $\zeta(2)$ et $\zeta(3)$, Ast\'erisque 61 (1979) 11--13.
\end{thebibliography}
\end{document}
